\chapter{The Jacobian as a class group}
This chapter develops the class-group model of the Jacobian over an arbitrary field $K$
with $2 \in K^\times$; we write $f_K$ for the image of $f$ in $K[x]$.

\begin{definition}[Hyperelliptic curve]
\label{def:hypcurve}
\lean{HyperellipticCurve, HyperellipticCurve.genus, HyperellipticCurve.IsOdd,
  HyperellipticCurve.map, HyperellipticCurve.ofSquarefree, HyperellipticCurve.ofMonicOdd,
  Polynomial.discrOfDegree, Polynomial.discrOfDegree_add_two}
\leanok
Let $R$ be a commutative ring in which $2$ is invertible. A hyperelliptic curve of genus $g$
over $R$ is given by a polynomial $f \in R[x]$ of degree at most $2g + 2$ such that the
discriminant of the binary form $Z^{2g+2} f(X/Z)$ of degree $2g+2$ is a unit of $R$; we write
$C \colon y^2 = f(x)$. When $R \ne 0$, the degree of $f$ is $2g+1$ or $2g+2$, and in the first
case the leading coefficient of $f$ is a unit. The model is of odd degree if $f$ is monic of
degree $2g+1$; the curve then has exactly one point at infinity. Every ring homomorphism
$\varphi \colon R \to S$ (with $2$ invertible in $S$) maps $C$ to the curve
$y^2 = f^\varphi(x)$ of the same genus over $S$, and a nonconstant squarefree polynomial over a
perfect field defines a curve.
\end{definition}

\begin{remark}[Formalization]
The discriminant of the binary form is \texttt{Polynomial.discrOfDegree}: Mathlib's
discriminant matrix with the degree replaced by $2g+2$, so that it is a polynomial in the
coefficients of $f$ and commutes with every ring homomorphism. It equals the discriminant of $f$
when $\deg f = 2g+2$ and $\operatorname{lc}(f)^2$ times it when $\deg f = 2g+1$; the general
form of the latter is the drop-by-one formula $D_{n+2}(f) = a_{n+1}^2\, D_{n+1}(f)$ for
$\deg f \le n+1$ (\texttt{Polynomial.discrOfDegree\_add\_two}), proved by expanding Mathlib's
modified Sylvester matrix along its normalized last row. The invertibility of $2$ is a parameter
of the type (\texttt{[Invertible (2 : R)]}); the genus is part of the data, so base change needs
no hypothesis. The coordinate ring, function field and Jacobian below are still indexed by the
polynomial $f$; a curve supplies their hypotheses (odd degree, squarefreeness and the
non-square property over a field) through instances.
\end{remark}

\begin{definition}[Coordinate ring and function field]
\label{def:hypring}
\lean{ChabautyColeman.HypRing, ChabautyColeman.HypField,
  ChabautyColeman.HypRing.iota, ChabautyColeman.HypField.iota}
\leanok
$\A{K} := K[x][Y]/(Y^2 - f_K)$, written $K[x,y]$, and
$\F{K} := K(x)[Y]/(Y^2 - f_K)$. The hyperelliptic involution $\iota$ is the $K[x]$-
(resp.\ $K(x)$-) algebra automorphism with $\iota(y) = -y$.
\end{definition}

\begin{lemma}[Irreducibility]
\label{lem:hyp-irred}
\lean{ChabautyColeman.irreducible_Ysq_sub_f, ChabautyColeman.not_isSquare_algebraMap_ratFunc,
  ChabautyColeman.instFactNotIsSquareOfOdd, IsIntegrallyClosed.isSquare_algebraMap_iff,
  HyperellipticCurve.instFactNotIsSquareF, Squarefree.isUnit_of_isSquare,
  ChabautyColeman.instFactHypField, QuadraticAlgebra.finrank_eq_two,
  ChabautyColeman.HypRing.toHypField, ChabautyColeman.HypRing.toHypField_injective,
  ChabautyColeman.HypRing.instIsDomain}
\leanok
\uses{def:hypring}
For every field $K$, the polynomial $Y^2 - f_K$ is irreducible over $K(x)$; hence $\F{K}$
is a field, quadratic over $K(x)$, and $\A{K}$ is a domain, as it embeds into $\F{K}$.
(The extension is moreover Galois with group $\{1,\iota\}$; this refinement is not used
below.)
\end{lemma}
\begin{proof}
\leanok
A root would give $f_K = r^2$ with $0 \ne r \in K(x)$; comparing degrees,
$\deg f = 2\deg r$ is even, contradicting $\deg f = 2g+1$. \emph{Formalization.} The
statements assume only that~$f_K$ is not a square in~$K[x]$: as $K[x]$ is integrally closed,
$f_K$ is then not a square in~$K(x)$ either. One instance derives this hypothesis from the odd
degree, by the parity argument on \texttt{RatFunc.intDegree}, and another one from a curve over
a field, whose polynomial is squarefree of positive degree.
\end{proof}

\begin{lemma}[Norm form]
\label{lem:normform}
\lean{ChabautyColeman.HypField.norm_eq, ChabautyColeman.HypField.trace_eq,
  ChabautyColeman.HypField.mul_star_eq_norm, ChabautyColeman.HypRing.norm_eq,
  ChabautyColeman.HypRing.trace_eq, ChabautyColeman.HypRing.mul_star_eq_norm}
\leanok
\uses{lem:hyp-irred}
For $\alpha = a + by \in \F{K}$ one has
$\Nm_{\F{K}/K(x)}(\alpha) = \alpha\,\iota(\alpha) = a^2 - b^2 f_K$
and $\operatorname{Tr}(\alpha) = 2a$. The same formulas hold for the norm and trace of
$\A{K}$ over $K[x]$ (and over any base ring).
\end{lemma}
\begin{proof}
\leanok
Computation in the basis $(1,y)$.
\end{proof}

\begin{lemma}[Integral elements]
\label{lem:integral-char}
\lean{ChabautyColeman.HypField.isIntegral_iff, ChabautyColeman.HypField.isIntegral_y,
  ChabautyColeman.HypRing.isIntegral_toHypField,
  RatFunc.exists_algebraMap_eq_of_sq_mul_squarefree}
\leanok
\uses{lem:hyp-irred, lem:normform}
If $f_K$ is squarefree (and $2 \in K^\times$), then $a + by \in \F{K}$ is integral over
$K[x]$ if and only if $a, b \in K[x]$, i.e.\ if and only if $a + by$ lies in the image of
$\A{K}$.
\end{lemma}
\begin{proof}
\leanok
($\Leftarrow$) $y$ satisfies $Y^2 = f_K$. ($\Rightarrow$) $\alpha$ and $\iota\alpha$ are
integral, so $2a \in K[x]$ and $a^2 - b^2 f_K \in K[x]$ ($K[x]$ being integrally closed);
hence $b^2 f_K \in K[x]$ forces $b \in K[x]$: writing $b = p/q$ in lowest terms with $q$
monic, $p^2 f_K = c\,q^2$ and $\gcd(p,q) = 1$ give $q^2 \mid f_K$, so $q = 1$ by
squarefreeness. (This lowest-terms argument replaces a per-prime valuation computation.)
\end{proof}

\begin{theorem}[The coordinate ring is Dedekind]
\label{thm:hypring-dedekind}
\lean{ChabautyColeman.HypRing.instIsIntegralClosure,
  ChabautyColeman.HypRing.instIsDedekindDomain, ChabautyColeman.HypField.instIsSeparable,
  Algebra.IsSeparable.of_finrank_le_two}
\leanok
\uses{lem:integral-char}
If $f_K$ is squarefree, then $\A{K}$ is the integral closure of $K[x]$ in $\F{K}$ and is a
Dedekind domain, free over $K[x]$ with basis $(1,y)$.
\end{theorem}
\begin{proof}
\leanok
\uses{lem:integral-char, lem:hyp-irred}
By Lemma~\ref{lem:integral-char} the integral closure is $K[x] \oplus K[x]\,y = \A{K}$;
the integral closure of the principal ideal domain $K[x]$ in the finite separable
extension $\F{K}/K(x)$ is Dedekind. Separability holds because the extension is quadratic
and $2 \in K^\times$: any extension of degree $\le 2$ in characteristic $\ne 2$ is
separable. Freeness with basis $(1,y)$ is immediate from the construction.
\emph{Formalization.} The Dedekind property is Mathlib's
\texttt{IsIntegralClosure.isDedekindDomain}.
\end{proof}

\begin{lemma}[Units]
\label{lem:hypring-units}
\lean{ChabautyColeman.HypRing.isUnit_iff}
\leanok
\uses{lem:normform}
$(\A{K})^\times = K^\times$ for every field $K$. Only the odd degree of $f$ is used; no
assumption on the characteristic or on squarefreeness is needed.
\end{lemma}
\begin{proof}
\leanok
$\Nm(u) = a^2 - b^2 f_K \in K[x]^\times = K^\times$; if $b \ne 0$, the right-hand side of
$a^2 = c + b^2 f_K$ has odd degree $2\deg b + 2g + 1$, while the left-hand side has even
degree. So $b = 0$ and $a$ is constant.
\end{proof}

\begin{definition}[Jacobian and Abel--Jacobi map]
\label{def:hypjac}
\lean{ChabautyColeman.HypJac, ChabautyColeman.aj, ChabautyColeman.pointIdeal}
\leanok
\uses{thm:hypring-dedekind}
$\J{K} := \Cl(\A{K})$. For an affine point $P = (x_0,y_0) \in C(K)$ set
$\aj(P) := [\mathfrak m_P] \in \J{K}$, where $\mathfrak m_P = (x - x_0,\, y - y_0)$ is the
maximal ideal of $P$, and $\aj(\infty) := 0$. More generally, the class of a $K$-rational
divisor supported on affine points is the class of the corresponding fractional ideal.
\end{definition}

\begin{remark}[Formalization]
In Lean, $\Cl$ is Mathlib's multiplicative \texttt{ClassGroup}, so the identity class is
written $1$; classes of general ideals are obtained through \texttt{ClassGroup.mk0} and
need no separate definition.
\end{remark}

\begin{lemma}[The involution acts as inversion on point classes]
\label{lem:aj-inverse}
\lean{ChabautyColeman.aj_inv, ChabautyColeman.pointIdeal_mul_pointIdeal_neg}
\leanok
\uses{def:hypjac, lem:mumford-norm}
$\divisor(x - x_0)_{\mathrm{aff}} = \mathfrak m_P \cdot \mathfrak m_{\iota P}$, hence
$\aj(\iota P) = -\aj(P)$. (That $\iota$ inverts \emph{all} of $\J{K}$ is
Lemma~\ref{lem:iota-norm}.)
\end{lemma}
\begin{proof}
\leanok
This is the special case $u = x - x_0$, $v = y_0$ of the Mumford norm identity
(Lemma~\ref{lem:mumford-norm}), since classes of principal ideals vanish.
\end{proof}

\begin{lemma}[Conjugate ideal and norm]
\label{lem:iota-norm}
\lean{ChabautyColeman.mul_map_star_eq_map_relNorm, ChabautyColeman.classGroup_mk0_map_star,
  ChabautyColeman.classGroup_mk0_mumfordIdeal_map_star, Squarefree.dvd_of_sq_dvd_mul_sq}
\leanok
\uses{thm:hypring-dedekind, lem:normform, lem:integral-char}
For every ideal $\mathfrak a$ of $\A{K}$:
$\mathfrak a \cdot \iota(\mathfrak a) = (\Nm \mathfrak a)\,\A{K}$, where
$\Nm \mathfrak a \subseteq K[x]$ is the relative ideal norm. Hence $\iota$ induces
inversion on $\J{K}$.
\end{lemma}
\begin{proof}
\leanok
Since $K[x]$ is a principal ideal domain, $\Nm \mathfrak a = (n_0)$ is principal, and no
factorization into primes is needed: the inclusion $\supseteq$ holds because
$N(a) = a\,\iota(a) \in \mathfrak a\,\iota(\mathfrak a)$ for $a \in \mathfrak a$. For
$\subseteq$, take $t = a\,\iota(b)$ with $a, b \in \mathfrak a$; then
$t + \iota t = N(a+b) - N(a) - N(b) \in (n_0)$ gives $n_0 \mid 2\,t_{\mathrm{re}}$, and
$t\,\iota t = N(a)N(b) \in (n_0^2)$ gives
$n_0^2 \mid t_{\mathrm{re}}^2 - f\,t_{\mathrm{im}}^2$; since $2 \in K^\times$ and $f$ is
squarefree, $n_0$ divides both components of $t$ (using that $n^2 \mid f b^2$ with $f$
squarefree forces $n \mid b$ in a unique factorization domain).
\end{proof}

\begin{remark}[Formalization]
Lemma~\ref{lem:iota-norm} is formalized in
\texttt{HyperellipticCurves.Archive.ConjugateIdealNorm}; the main theorems do not depend on it,
since both places that invert a class invert a \emph{Mumford} ideal, where the Mumford norm
identity of Lemma~\ref{lem:mumford-norm} gives the inversion directly.
\end{remark}

\begin{definition}[Mumford ideals]
\label{def:mumford-ideal}
\lean{ChabautyColeman.mumfordIdeal, ChabautyColeman.mumfordHom,
  ChabautyColeman.mumfordQuotEquiv, ChabautyColeman.ker_mumfordHom}
\leanok
\uses{def:hypjac}
For $u, v \in K[x]$, the \emph{Mumford ideal} is
$\mathfrak a_{u,v} := (u,\, y - v) \subseteq \A{K}$; point ideals are the case
$(x - x_0,\, y - y_0)$. If $u \mid v^2 - f_K$, then the evaluation
$p + qy \mapsto p + qv \bmod u$ is a surjective ring homomorphism
$\A{K} \to K[x]/(u)$ with kernel $\mathfrak a_{u,v}$, so
$\A{K}/\mathfrak a_{u,v} \cong K[x]/(u)$; in particular $\mathfrak a_{u,v}$ is a
proper nonzero ideal for nonconstant $u$.
\end{definition}

\begin{lemma}[Mumford norm identity]
\label{lem:mumford-norm}
\lean{ChabautyColeman.mumfordIdeal_mul_map_star, ChabautyColeman.map_star_mumfordIdeal}
\leanok
\uses{def:mumford-ideal}
If $u \mid v^2 - f_K$ with $f_K$ squarefree (and $2 \in K^\times$), then
$\iota(\mathfrak a_{u,v}) = \mathfrak a_{u,-v}$ and
$\mathfrak a_{u,v} \cdot \iota(\mathfrak a_{u,v}) = (u)$.
\end{lemma}
\begin{proof}
\leanok
$\subseteq$: all four generator products are divisible by $u$, using
$(y-v)(y+v) = f_K - v^2 = -uw$ with $w := (v^2-f_K)/u$. $\supseteq$: a common prime $\pi$
of $(u, 2v, w)$ would give $\pi \mid v$ and $\pi^2 \mid uw = v^2 - f_K$, so
$\pi^2 \mid f_K$, contradicting squarefreeness; a B\'ezout identity then exhibits $u$ as
a combination of $u^2$, $u\cdot 2v$ and $uw$, all of which lie in the product.
\end{proof}

\begin{lemma}[The divisor of $y - v$]
\label{lem:mumford-divisor-y}
\lean{ChabautyColeman.mumfordIdeal_mul_cofactor,
  ChabautyColeman.classGroup_mk0_mumfordIdeal_cofactor}
\leanok
\uses{def:mumford-ideal, def:hypjac}
If $v^2 - f_K = u\,w$ (with $f_K$ squarefree, $2 \in K^\times$), then
$\mathfrak a_{u,v} \cdot \mathfrak a_{w,v} = (y - v)$ is principal; hence
$[\mathfrak a_{w,v}] = [\mathfrak a_{u,v}]^{-1}$ in $\J{K}$. Combined with
Lemma~\ref{lem:mumford-norm} this yields Cantor's reduction step:
$[\mathfrak a_{u,v}] = [\mathfrak a_{w,-v}]$ with
$\deg w \le \max(2\deg v, 2g+1) - \deg u$.
\end{lemma}
\begin{proof}
\leanok
$\subseteq$: the generator products are divisible by $y - v$, using
$uw = v^2 - f_K = -(y-v)(y+v)$. $\supseteq$: the same B\'ezout identity for
$(u, 2v, w)$ as in Lemma~\ref{lem:mumford-norm} exhibits $y - v$ in the product,
using $(y-v)\cdot 2v = -uw - (y-v)^2$.
\end{proof}

\begin{lemma}[Semi-reduction]
\label{lem:mumford-semireduced}
\lean{ChabautyColeman.exists_span_singleton_mul_mumfordIdeal,
  ChabautyColeman.exists_mk0_eq_mk0_mumfordIdeal}
\leanok
\uses{def:mumford-ideal, lem:normform, def:hypjac}
Every nonzero ideal $\mathfrak a$ of $\A{K}$ has the form
$\mathfrak a = (c)\,\mathfrak a_{u,v}$ with $c, u, v \in K[x]$, $c, u \ne 0$ and
$u \mid v^2 - f_K$; hence every class in $\J{K}$ is the class of a Mumford ideal.
\end{lemma}
\begin{proof}
\leanok
Hermite normal form over the principal ideal domain $K[x]$, with no factorization into
primes: let $c$ generate the image of $\mathfrak a$ under the coefficient projection
$a + by \mapsto b$ and let $a_0$ generate $\mathfrak a \cap K[x]$; multiplying elements
by $y$ shows $c \mid a_0$ and that $c$ divides the $K[x]$-part of a lift $z_0$ of $c$,
giving $\mathfrak a = (a_0, z_0) = (c)(u, y - v)$ with $a_0 = cu$, $z_0 = c(y-v)$.
Stability of $\mathfrak a$ under multiplication by $y$, applied to $z_0$, yields
$u \mid v^2 - f_K$. (Nonvanishing of $a_0$ uses the norm form:
$0 \ne \Nm(z) \in \mathfrak a \cap K[x]$.)
\end{proof}

\begin{lemma}[Existence of reduced representatives]
\label{lem:mumford-reduce}
\lean{ChabautyColeman.exists_mk0_mumfordIdeal_reduced}
\leanok
\uses{lem:mumford-semireduced, lem:mumford-divisor-y, lem:mumford-norm}
Every class in $\J{K}$ is the class of a \emph{reduced} Mumford ideal
$\mathfrak a_{u,v}$: $u$ monic, $\deg v < \deg u \le g$, $u \mid v^2 - f_K$.
\end{lemma}
\begin{proof}
\leanok
By Lemma~\ref{lem:mumford-semireduced} the class contains a Mumford ideal. The Mumford
ideal is unchanged by unit rescaling of $u$ and by replacing $v$ with $v \bmod u$;
while $\deg u > g$, the class-preserving replacement
$[\mathfrak a_{u,v}] = [\mathfrak a_{w,-v}]$ with $v^2 - f_K = uw$
(Lemmas~\ref{lem:mumford-divisor-y} and~\ref{lem:mumford-norm}) strictly decreases
$\deg u$: after reducing $v$ modulo $u$,
$\deg w \le \max(2\deg v, 2g+1) - \deg u < \deg u$. Conclude by induction on $\deg u$.
\end{proof}

\begin{lemma}[Uniqueness of reduced representatives]
\label{lem:mumford-exists}
\lean{ChabautyColeman.mumfordIdeal_reduced_unique, ChabautyColeman.mem_mumfordIdeal_iff,
  ChabautyColeman.comap_mumfordIdeal}
\leanok
\uses{thm:hypring-dedekind, lem:mumford-norm, def:mumford-ideal,
  lem:mumford-norm, lem:mumford-divisor-y, lem:mumford-reduce}
The reduced Mumford ideal representing a class (Lemma~\ref{lem:mumford-reduce}) is
unique. (Cantor's composition-and-reduction algorithm computes the reduced
representative of a product of two reduced classes; existence is
Lemma~\ref{lem:mumford-reduce}.)
\end{lemma}
\begin{proof}
\leanok
Let $\mathfrak a_1, \mathfrak a_2$ be reduced Mumford ideals in the same class. Then
$\mathfrak a_1\,\iota(\mathfrak a_2)$ is an integral principal ideal $(h)$, and
$(\Nm h) = (u_1 u_2)$ by the norm identity, of degree $\le 2g$. Since $\Nm(a + by) =
a^2 - b^2 f_K$ has degree $\ge 2g+1$ whenever $b \ne 0$ (parity), $h = a \in K[x]$.
Multiplying $(h) = \mathfrak a_1 \iota(\mathfrak a_2)$ and its conjugate by
$\mathfrak a_2$ resp.\ $\mathfrak a_1$ and taking traces to $K[x]$ (the trace of
$\mathfrak a_{u,v}$ is $(u)$, immediate from the quotient description) gives
$u_1 \sim a \sim u_2$, hence $u_1 = u_2$ by monicity; cancelling the principal factor
yields $\mathfrak a_1 = \mathfrak a_2$, and then $u \mid v_2 - v_1$ with
$\deg(v_2 - v_1) < \deg u$ forces $v_1 = v_2$. (The norm-degree parity argument
replaces the usual pole count at the place at infinity.)
\end{proof}

\begin{definition}[Base change]
\label{def:basechange}
\lean{ChabautyColeman.HypJac.map, ChabautyColeman.HypRing.baseChange,
  ChabautyColeman.HypJac.map_aj, ChabautyColeman.pointIdeal_map}
\leanok
\uses{def:hypjac}
For a field extension $K \subseteq K'$, the homomorphism $\J{K} \to \J{K'}$ induced by
ideal extension along $\A{K} \to \A{K'}$. It is compatible with the Abel--Jacobi maps.
\emph{Formalization:} this is Mathlib's \texttt{ClassGroup.extendedHom} applied to the
base change of coordinate rings.
\end{definition}

\begin{lemma}[Injectivity criterion]
\label{lem:basechange-crit}
\lean{ChabautyColeman.hypJac_map_injective_iff}
\leanok
\uses{def:basechange, def:hypjac}
$\J{K} \to \J{K'}$ is injective if and only if every nonzero ideal of $\A{K}$ whose
extension to $\A{K'}$ is principal is itself principal.
\end{lemma}
\begin{proof}
\leanok
Kernel triviality for the group homomorphism, using the surjectivity of
$\mathfrak a \mapsto [\mathfrak a]$ and the compatibility of the base-change map with
ideal extension.
\end{proof}

\begin{lemma}[Contraction]
\label{lem:basechange-contract}
\lean{ChabautyColeman.comap_map_baseChange}
\leanok
\uses{def:basechange}
For every ideal $\mathfrak a$ of $\A{K}$ and every field extension $K \subseteq K'$:
$(\mathfrak a \A{K'}) \cap \A{K} = \mathfrak a$.
\end{lemma}
\begin{proof}
\leanok
Choose a $K$-linear splitting $p \colon K' \to K$ of the inclusion and apply it to all
polynomial coefficients; the resulting additive retraction
$\pi \colon \A{K'} \to \A{K}$ satisfies $\pi(\iota(a)w) = a\,\pi(w)$, and applying
it to an element of $\mathfrak a \A{K'} \cap \A{K}$, written as a combination
$\sum \iota(a_j) w_j$, lands in $\mathfrak a$. No flatness theory is needed.
\end{proof}

\begin{lemma}[Galois action and the cocycle]
\label{lem:basechange-galact}
\lean{ChabautyColeman.HypRing.galAct, ChabautyColeman.HypRing.galAct_algebraMap,
  ChabautyColeman.HypRing.exists_unit_galAct_eq}
\leanok
\uses{def:basechange, lem:hypring-units}
A $K$-automorphism $\sigma$ of $K'$ acts on $\A{K'}$ coefficientwise, fixing the
image of $\A{K}$ and stabilizing extended ideals. If
$\mathfrak a\,\A{K'} = (h)$ is principal, then $\sigma(h) = c_\sigma\, h$ for a
constant $c_\sigma \in (K')^\times$.
\end{lemma}
\begin{proof}
\leanok
$(\sigma(h)) = \sigma(\mathfrak a \A{K'}) = \mathfrak a \A{K'} = (h)$, so
$\sigma(h)$ and $h$ are associated; the unit is a constant by
Lemma~\ref{lem:hypring-units}.
\end{proof}

\begin{lemma}[Fixed-point descent]
\label{lem:basechange-fixed}
\lean{ChabautyColeman.HypRing.exists_algebraMap_eq_of_forall_galAct}
\leanok
\uses{lem:basechange-galact}
For a finite Galois extension $K \subseteq K'$, an element of $\A{K'}$ fixed by the
full Galois action lies in the image of $\A{K}$.
\end{lemma}
\begin{proof}
\leanok
Coefficientwise from $(K')^{\mathrm{Gal}(K'/K)} = K$.
\end{proof}

\begin{lemma}[Descent of principality along Galois extensions]
\label{lem:basechange-galois}
\lean{ChabautyColeman.HypRing.isPrincipal_of_map_isPrincipal,
  ChabautyColeman.hypJac_map_injective_of_isGalois}
\leanok
\uses{lem:basechange-crit, lem:basechange-contract, lem:basechange-galact,
  lem:basechange-fixed}
Let $K \subseteq K'$ be a finite Galois extension. If an ideal $\mathfrak a$ of $\A{K}$
extends to a principal ideal of $\A{K'}$, then $\mathfrak a$ is principal. Consequently
(by Lemma~\ref{lem:basechange-crit}) the map $\J{K} \to \J{K'}$ is injective.
\end{lemma}
\begin{proof}
\leanok
Let $\mathfrak a\,\A{K'} = (h)$; then $h \ne 0$ unless $\mathfrak a = 0$ (which is
trivially principal). Lemma~\ref{lem:basechange-galact} gives
$c \colon \mathrm{Gal}(K'/K) \to (K')^\times$ with $\sigma(h) = c_\sigma h$; cancelling
$h$ from the two evaluations of $(\sigma\tau)(h)$ shows $c$ is a $1$-cocycle. By
Hilbert's Theorem~90 it is a coboundary, $c_\sigma = \sigma(\beta)/\beta$ for some
$\beta \in (K')^\times$. Then $h' = \beta^{-1} h$ is a Galois-fixed generator of
$\mathfrak a\,\A{K'}$, so $h'$ comes from some $w \in \A{K}$ by
Lemma~\ref{lem:basechange-fixed}. Contracting via Lemma~\ref{lem:basechange-contract}
gives $\mathfrak a = \mathfrak a\,\A{K'} \cap \A{K} = w\A{K'} \cap \A{K} = w\A{K}$.
\end{proof}

\begin{theorem}[Base-change injectivity]
\label{thm:basechange-inj}
\lean{ChabautyColeman.hypJac_map_injective, ChabautyColeman.mumfordIdeal_map}
\leanok
\uses{def:basechange, lem:mumford-reduce, lem:mumford-exists}
For \emph{every} field extension $K \subseteq K'$ (with the standing hypotheses on $f$
over both fields), the map $\J{K} \to \J{K'}$ is injective. In particular this applies
to $K' \in \{\QQ_p, \Qq, \Cp\}$ over $K = \QQ$; no finiteness, separability, or
characteristic assumption on the extension is needed.
\end{theorem}
\begin{proof}
\leanok
\uses{lem:mumford-reduce, lem:mumford-exists}
Represent two classes with equal images by their reduced Mumford representatives
$(u_1, v_1)$ and $(u_2, v_2)$ (Lemma~\ref{lem:mumford-reduce}). Base change sends
Mumford ideals to Mumford ideals and preserves the reducedness conditions (monicity,
the degree bounds, $u \mid v^2 - f$), so the base-changed pairs are reduced
representatives of the common image class, and by uniqueness
(Lemma~\ref{lem:mumford-exists}) they coincide over $K'$. Since all coefficients lie in
$K$ and $K \to K'$ is injective, $(u_1, v_1) = (u_2, v_2)$, so the classes agree.
\end{proof}

\begin{remark}[Formalization]
Base-change injectivity and its Galois-descent chain
(Lemmas~\ref{lem:basechange-crit}--\ref{lem:basechange-galois},
Theorem~\ref{thm:basechange-inj}) are formalized in
\texttt{HyperellipticCurves.Archive.BaseChangeInjectivity}; the main theorems do not depend on
them, since the main development only uses the base-change \emph{homomorphism} of
Definition~\ref{def:basechange}. The comparison of reduced representatives in the proof of
Theorem~\ref{thm:basechange-inj} needs no descent and works for any field extension, whereas
descent through finite Galois subextensions (Hilbert~90) covers only algebraic extensions: a
transcendental extension such as $\QQ_p/\QQ$ would require an additional specialization
argument.
\end{remark}
